% ================================================================
%  §9  Constrained Optimization Algorithms
% ================================================================
\section{Constrained Optimization Algorithms}\label{sec:constrained-algorithms}

The previous section established \emph{what} characterizes a constrained optimum --- KKT conditions, duality, second-order tests.
This section addresses \emph{how} to find one.
The algorithmic landscape for constrained optimization is richer than the unconstrained case, because the feasible region $X$ introduces geometric structure that each method exploits differently.

The organizing principle: the structure of $X$ determines the algorithm.
Linear inequality constraints lead to active-set and simplex-type methods; simple constraints (boxes, balls, simplices) favor projected gradient and Frank--Wolfe; general convex constraints are handled by barrier and interior point methods; and the dual perspective transforms complex constraints into simple ones at the cost of a nonsmooth objective.


% ────────────────────────────────────────────────────────────────
\subsection{Equality-Constrained Quadratic Programs}\label{ssec:eq-constrained-qp}

These serve as the foundation --- every more complex method eventually reduces to solving a sequence of equality-constrained QPs.

The problem is $\min\{\frac{1}{2} x^T Q x + q^T x : Ax = b\}$ with $A \in \R^{m \times n}$, $\rank(A) = m < n$.
The KKT conditions yield the symmetric indefinite system
\begin{equation}\label{eq:kkt-system}
	\begin{bmatrix} Q & A^T \\ A & 0 \end{bmatrix}
	\begin{bmatrix} x \\ \mu \end{bmatrix}
	= \begin{bmatrix} -q \\ b \end{bmatrix}.
\end{equation}
Three solution strategies exist, each with its own trade-offs.

\paragraph{Schur complement (reduced KKT)}

When $Q$ is nonsingular, eliminate $x$ to get the smaller system $[AQ^{-1}A^T]\mu = -b - AQ^{-1}q$.
The Schur complement $M = AQ^{-1}A^T \in \R^{m \times m}$ is positive semidefinite when $Q \succeq 0$ and smaller than the original system, but can be dense even when $A$ and $Q$ are sparse.
Preconditioned conjugate gradient on $M$ avoids forming it explicitly.

\paragraph{Null space method}

Partition $A = [A_B \mid A_N]$ with $A_B$ invertible.
The constraint $Ax = b$ gives $x_B = A_B^{-1}(b - A_N x_N)$, so $x = Dx_N + d$ where $D = [-A_B^{-1}A_N;\; I]^T$.
The \textit{reduced Hessian} $\tilde{Q} = D^T Q D \in \R^{(n-m) \times (n-m)}$ governs the reduced problem $\min_{x_N} \frac{1}{2} x_N^T \tilde{Q}\, x_N + (D^T Q d + D^T q)^T x_N$.
The method generalizes to any basis of $\ker(A)$ and handles singular $Q$.

\paragraph{Cholesky on the KKT system}

For moderate $n$, direct factorization of the KKT system at $\mathcal{O}(n^3)$ is often the best choice, providing both $x$ and $\mu$ simultaneously.


% ────────────────────────────────────────────────────────────────
\subsection{Active Set Methods}\label{ssec:active-set}

Active set methods solve inequality-constrained QPs by maintaining an estimate of the active constraints and solving a sequence of equality-constrained subproblems.

The key insight: if we knew $\mathcal{A}(x^*) = \{i : A_i x^* = b_i\}$, the problem would reduce to linear algebra.
Since we don't, we estimate it iteratively, using dual information to guide updates.

\begin{algorithm}
	\caption{Active Set Method for Quadratic Programs (ASMQP)}
	\label{alg:ASMQP}
	\begin{algorithmic}[1]
		\Require $Q \in \R^{n \times n}$, $q \in \R^n$, $A \in \R^{m \times n}$, $b \in \R^m$, feasible $x^0$
		\Ensure Optimal solution to $(P)\; \min\{\frac{1}{2} x^T Q x + q^T x : Ax \leq b\}$
		\Procedure{ASMQP}{$Q, q, A, b, x^0$}
			\State $x \gets x^0$;\; $B \gets \mathcal{A}(x)$
			\While{true}
				\State Solve $(P_B)\; \min\{f(x) : A_B x = b_B\}$; let $(\bar{x}, \bar{\mu}_B)$ be the KKT pair
				\If{$A_i \bar{x} \leq b_i$ for all $i \notin B$} \Comment{Candidate feasible?}
					\If{$\bar{\mu}_B \geq 0$} \Return $\bar{x}$ \Comment{KKT satisfied: optimal}
					\Else
						\State Remove $h = \argmin_{i \in B}\{\bar{\mu}_i\}$ from $B$ \Comment{Drop most violated dual}
					\EndIf
				\Else
					\State $d \gets \bar{x} - x$;\;
						$\bar{\alpha} \gets \min\{(b_i - A_i x)/(A_i d) : A_i d > 0,\, i \notin B\}$
					\State $x \gets x + \bar{\alpha}\, d$;\; $B \gets \mathcal{A}(x)$
				\EndIf
			\EndWhile
		\EndProcedure
	\end{algorithmic}
\end{algorithm}

Each iteration either moves to a strictly better feasible point or improves the dual solution.
Since there are at most $2^m$ possible active sets and the objective decreases, the method terminates in finite time.
Linear independence of $A_B$ is required for well-posedness.

\paragraph{Box-constrained problems}

For $\min\{\frac{1}{2} x^T Q x + q^T x : \underline{x} \leq x \leq \overline{x}\}$, the active set structure is particularly elegant.
Each active constraint corresponds to a variable fixed at its bound.
The reduced problem involves only the ``free'' variables: $\min\{\frac{1}{2} x_F^T Q_{FF} x_F + \tilde{q}_F^T x_F\}$ where $\tilde{q}$ absorbs the contribution of fixed variables.
Linear independence is automatic, initialization is trivial (any point in the box), and active set updates add or remove single bounds.

Active set methods extend to general nonlinear objectives by replacing the equality subproblem with an unconstrained optimization on the reduced space.


% ────────────────────────────────────────────────────────────────
\subsection{Projected Gradient Methods}\label{ssec:projected-gradient}

Projected gradient methods maintain feasibility at every iteration while making progress toward optimality.
When $-\nabla f(x)$ is infeasible (it points outside $X$), we project it onto the cone of feasible directions.

\paragraph{Box-constrained projected gradient}

For box constraints, the projection decomposes into independent scalar projections: simply zero out gradient components that would push a variable past its bound.
The step size is bounded by the distance to the nearest violated bound.

\begin{algorithm}
	\caption{Box-Constrained Projected Gradient (BCPGM)}
	\label{alg:BCPGM}
	\begin{algorithmic}[1]
		\Require $f$, bounds $\underline{x}, \overline{x}$, initial feasible $x$, tolerance $\varepsilon$
		\Procedure{BCPGM}{$f, \underline{x}, \overline{x}, x, \varepsilon$}
			\While{true}
				\State $d \gets -\nabla f(x)$;\; $\bar{\alpha} \gets \infty$
				\For{$i = 1, \dots, n$}
					\If{$d_i < 0$ and $x_i = \underline{x}_i$, or $d_i > 0$ and $x_i = \overline{x}_i$}
						\State $d_i \gets 0$ \Comment{Project out infeasible component}
					\ElsIf{$d_i \neq 0$}
						\State $\bar{\alpha} \gets \min(\bar{\alpha},\; \text{distance to bound along } d_i)$
					\EndIf
				\EndFor
				\If{$\inner{\nabla f(x), d} \geq -\varepsilon$} \Return $x$ \EndIf
				\State $\alpha \gets \text{line search}(f, x, d, \bar{\alpha})$;\; $x \gets x + \alpha d$
			\EndWhile
		\EndProcedure
	\end{algorithmic}
\end{algorithm}

\paragraph{Goldstein's method}

An alternative: take the gradient step first, then project onto $X$:
$y^k = x^k - \alpha^k \nabla f(x^k)$, then $x^{k+1} = p_X(y^k)$.
This is not a descent method, but works well when projection onto $X$ is cheap.
For $L$-smooth convex functions with $\alpha = 1/L$: $\mathcal{O}(LD^2/\varepsilon)$ iterations; with $\tau$-strong convexity: $\mathcal{O}((L/\tau)\log(1/\varepsilon))$ linear convergence.

Projection cost determines practical efficiency: $\mathcal{O}(n)$ for boxes ($p_X(y)_i = \max(\underline{x}_i, \min(y_i, \overline{x}_i))$), $\mathcal{O}(n)$ for balls, closed-form for simplices.

\paragraph{Rosen's method}

For general linear constraints $Ax \leq b$, project onto the tangent space of active constraints.
With $\bar{A} = A_{\mathcal{A}(x)}$ of full row rank:
$d = (I - \bar{A}^T[\bar{A}\bar{A}^T]^{-1}\bar{A})(-\nabla f(x))$.
If the projected gradient is zero, check dual feasibility (KKT multipliers $\lambda = -[\bar{A}\bar{A}^T]^{-1}\bar{A}\nabla f(x)$); if some $\lambda_h < 0$, drop constraint $h$ and re-project.


% ────────────────────────────────────────────────────────────────
\subsection{Frank--Wolfe (Conditional Gradient) Method}\label{ssec:frank-wolfe}

The Frank--Wolfe method replaces the projection step with a linear minimization over $X$.
At each iterate, it linearizes $f$ and solves $\min\{\inner{\nabla f(x), z} : z \in X\}$, then moves toward the minimizer.

\begin{algorithm}
	\caption{Frank--Wolfe Method (FW)}
	\label{alg:FW}
	\begin{algorithmic}[1]
		\Require $f$, feasible region $X$, initial $x \in X$, tolerance $\varepsilon$
		\Procedure{FW}{$f, X, x, \varepsilon$}
			\While{true}
				\State $\bar{x} \gets \argmin\{\inner{\nabla f(x), z} : z \in X\}$ \Comment{Linear subproblem}
				\State $d \gets \bar{x} - x$
				\If{$\inner{\nabla f(x), d} \geq -\varepsilon$} \Return $x$ \EndIf
				\State $\alpha \gets \text{line search}(f, x, d, 1)$ with $\alpha \leq 1$
				\State $x \gets x + \alpha d$
			\EndWhile
		\EndProcedure
	\end{algorithmic}
\end{algorithm}

The optimality condition $\inner{\nabla f(x), d} = 0$ implies $x$ is optimal (for convex $f$, the quantity $f(x) + \inner{\nabla f(x), d}$ provides a computable lower bound, giving an optimality certificate).

Each iteration requires solving one LP over $X$.
The method shines when $X$ has structure that makes linear optimization cheap: network flows, semidefinite constraints ($X = \{X : \tr(X) \leq 1, X \succeq 0\}$ reduces to a leading eigenvector computation), or norm balls (closed-form solutions).

\paragraph{Enhancements}

A starting feasible point can be found by solving an LP with a dummy objective.
Quadratic regularization (add $\frac{\gamma}{2}\norm{z - x}^2$) or trust region constraints ($\norm{z - x}_\infty \leq \tau$) improve convergence without complicating the subproblem much.
Away steps (move away from vertices, not just toward new ones) help avoid stalling at non-optimal extreme points.

\paragraph{Extensions}

For $f \notin C^1$: replace $\nabla f(x)$ with the cutting plane model $f_{\mathcal{B}}$, giving the constrained bundle method $\min\{f_{\mathcal{B}}(z) + \gamma\norm{z - \bar{x}}^2 : z \in X\}$.
For $f \in C^2$: incorporate second-order information via SQP (Sequential Quadratic Programming): $d^* = \argmin\{\inner{\nabla f(x), d} + \frac{1}{2} d^T \nabla^2 f(x)\, d : x + d \in X\}$, requiring a constrained QP per iteration.


% ────────────────────────────────────────────────────────────────
\subsection{Dual Methods}\label{ssec:dual-methods}

Dual methods work in the space of Lagrange multipliers, transforming complex primal constraints into simple dual constraints.

For $(P)\; \min\{\frac{1}{2} x^T Q x + q^T x : Ax \leq b\}$ with $Q \succ 0$, the dual function is
\begin{equation}\label{eq:dual-function-qp}
	\psi(\lambda) = \lambda^T b - \frac{1}{2}(A^T\lambda - q)^T Q^{-1}(A^T\lambda - q),
\end{equation}
continuously differentiable with $\nabla\psi(\lambda) = b - Ax(\lambda)$ where $x(\lambda) = Q^{-1}(A^T\lambda - q)$.
The dual problem $(D)\; \max\{\psi(\lambda) : \lambda \geq 0\}$ is a concave maximization with simple box constraints.

As $\lambda^k \to \lambda^*$, the primal recovery $x(\lambda^k) \to x^*$.
The dual provides valid lower bounds throughout: $\psi(\lambda^k) \leq \nu(P)$.
When $X$ has additional structure, a primal-feasible approximation $\hat{x}^k = p_X(x(\lambda^k))$ gives an upper bound, and the gap $f(\hat{x}^k) - \psi(\lambda^k)$ is a computable optimality certificate.

\paragraph{Advantages}

Step size issues (related to $\kappa = L/\tau$) are often less critical in the dual; simple constraints $\lambda \geq 0$ are easier than complex primal constraints; parallelization opportunities arise through decomposition.

\paragraph{Decomposition}

For problems with separable structure $(P)\; \min\{f(x) : Ax \leq b,\, Ex \leq d\}$, relaxing only the ``linking'' constraints $Ax \leq b$ gives a Lagrangian relaxation $\psi(\lambda) = \min\{f(x) + \lambda^T(b - Ax) : Ex \leq d\}$ that decomposes into smaller, independent subproblems --- enabling parallel solution.


% ────────────────────────────────────────────────────────────────
\subsection{Barrier Methods}\label{ssec:barrier}

Barrier methods transform the constrained problem into a sequence of unconstrained ones by penalizing proximity to the constraint boundary.

\begin{defn}[Logarithmic barrier]\label{def:barrier}
	For $X = \{x : Ax \leq b\}$, the barrier function is
	\begin{equation}\label{eq:log-barrier}
		f_\gamma(x) = f(x) - \gamma \sum_{i=1}^{m} \ln(b_i - A_i x),
	\end{equation}
	where $\gamma > 0$ is the barrier parameter.
	$f_\gamma \in C^2$ when $f \in C^2$; $f_\gamma \to +\infty$ at the boundary of $X$ (keeping iterates strictly interior); $f_\gamma$ is strictly convex when $f$ is convex.
\end{defn}

For each $\gamma > 0$, the unique minimizer $x_\gamma$ traces the \textit{central path} $\mathcal{C} = \{x_\gamma : \gamma \in (0, \infty)\}$.
As $\gamma \to 0$, $x_\gamma \to x^*$ (the constrained optimum); as $\gamma \to \infty$, $x_\gamma$ approaches the \textit{analytical center} $x_\infty = \argmax\{\prod_i (b_i - A_i x)\}$.

The \textit{path-following} strategy: start near the analytical center and decrease $\gamma$ toward zero, tracking the central path with Newton steps.
The logarithmic barrier is \textit{self-concordant}, ensuring that Newton's method behaves well.
If $x^k$ is close to $x_{\gamma^k}$, a few Newton steps yield $x^{k+1}$ close to $x_{\gamma^{k+1}}$ with $\gamma^{k+1} = \tau \gamma^k$ for $\tau < 1$.

Iteration complexity: $\mathcal{O}(\sqrt{m}\log(1/\varepsilon))$, independent of $n$ but depending on the number of constraints $m$.
Each Newton step costs $\mathcal{O}(n^3)$, so the total is $\mathcal{O}(n^3 \sqrt{m}\log(1/\varepsilon))$.


% ────────────────────────────────────────────────────────────────
\subsection{Primal-Dual Interior Point Methods}\label{ssec:ipm}

Primal-dual methods represent the most practically successful approach to barrier-based optimization, solving primal and dual problems simultaneously.

For the QP $\min\{\frac{1}{2} x^T Q x + q^T x : Ax \leq b\}$, introduce slacks $s \geq 0$ with $Ax + s = b$.
The KKT conditions are: primal feasibility ($Ax + s = b$, $s \geq 0$), dual feasibility ($Qx + q - A^T\lambda = 0$, $\lambda \geq 0$), and complementary slackness ($\lambda_i s_i = 0$).

The barrier method relaxes complementary slackness to $\lambda_i s_i = \gamma$ for all $i$, giving a smooth system solvable by Newton's method.
Using $\Lambda = \diag(\lambda)$ and $S = \diag(s)$, the Newton system for $(\Delta x, \Delta\lambda, \Delta s)$ is
\begin{equation}\label{eq:ipm-newton}
	\begin{bmatrix} Q & -A^T & 0 \\ A & 0 & I \\ 0 & S & \Lambda \end{bmatrix}
	\begin{bmatrix} \Delta x \\ \Delta\lambda \\ \Delta s \end{bmatrix}
	= \begin{bmatrix} r^D \\ r^P \\ \gamma u - \Lambda S u \end{bmatrix},
\end{equation}
where $r^P = b - Ax - s$ and $r^D = -(Qx + q - A^T\lambda)$ are residuals.

By eliminating $\Delta s$ and substituting, the system reduces to solving $(Q + A^T\Lambda S^{-1}A)\Delta x = \text{rhs}$.
The matrix $M = Q + A^T\Lambda S^{-1}A$ is positive definite when $Q \succeq 0$ (since $\Lambda S^{-1} \succ 0$), ensuring a unique solution.

\begin{algorithm}
	\caption{Primal-Dual Interior Point Method (IPMQP)}
	\label{alg:ipmqp}
	\begin{algorithmic}[1]
		\Require $Q, A, q, b$, tolerances $\varepsilon_P, \varepsilon_D, \varepsilon$, reduction factor $\rho$
		\Procedure{IPMQP}{$Q, A, q, b, \varepsilon_P, \varepsilon_D, \varepsilon, \rho$}
			\State Initialize $x$, $\lambda > 0$, $s > 0$
			\State $\gamma \gets \inner{\lambda, s} / m$
			\While{$\norm{r^P} > \varepsilon_P$ or $\norm{r^D} > \varepsilon_D$ or $\gamma m > \varepsilon$}
				\State Solve Newton system for $(\Delta x, \Delta\lambda, \Delta s)$
				\State $\alpha \gets 0.995 \cdot \max\{\beta : \lambda + \beta\Delta\lambda \geq 0,\; s + \beta\Delta s \geq 0\}$
				\State $(x, \lambda, s) \gets (x + \alpha\Delta x,\; \lambda + \alpha\Delta\lambda,\; s + \alpha\Delta s)$
				\State $\gamma \gets \rho \cdot \inner{\lambda, s} / m$
			\EndWhile
		\EndProcedure
	\end{algorithmic}
\end{algorithm}

\paragraph{Advanced techniques}

The \textit{predictor-corrector} approach solves the Newton system twice: first a ``prediction'' (affine scaling), then a ``correction'' that adds the nonlinear term $\Delta\Lambda\Delta S\, u$, reusing the matrix factorization.
Multiple centrality corrections further improve centering.
Adaptive selection of $\rho$ based on Newton step quality keeps the iterates near the central path.

\paragraph{Convergence}

The method achieves polynomial complexity $\mathcal{O}(\sqrt{m}\log(1/\varepsilon))$, with each iteration costing $\mathcal{O}(n^3)$ for the matrix factorization.
Superlinear convergence in the final phase is typical.
Primal and dual solutions converge simultaneously, providing both an optimal $x^*$ and optimal multipliers $\lambda^*$.
The framework extends directly to SOCP and SDP.
Numerical issues may arise in ill-conditioned or degenerate problems, where small slack or multiplier values cause instability in the $\Lambda S^{-1}$ scaling.